If $xyM\subseteq Q$ but $yM\not\subseteq Q$, then multiplication by $x$ kills a nonzero element of $M/Q$. Thus $x^nM\subseteq Q$ for some $n$, proving that $(Q:M)$ is primary. Put $\mathfrak{p}=r_M(Q)$.

If $Q_1,\ldots,Q_n$ are $\mathfrak{p}$-primary, then $Q=\bigcap_iQ_i$ has radical $\mathfrak{p}$ by \exref{4.20}. If $xm\in Q$ with $m\notin Q$, choose $i$ with $m\notin Q_i$. Then $x\in\mathfrak{p}$, so a sufficiently large power of $x$ sends $M$ into every $Q_i$, hence into $Q$. Thus $Q$ is $\mathfrak{p}$-primary.

For $m\in M$, write $(Q:m)=\{a\in A:am\in Q\}$. If $m\in Q$, this ideal is $A$. If $m\notin Q$, the map $A/(Q:m)\to M/Q$ induced by $a\mapsto am$ is injective. On any nonzero submodule of $M/Q$, elements outside $\mathfrak{p}$ act injectively and elements in $\mathfrak{p}$ act nilpotently. Hence $(Q:m)$ is $\mathfrak{p}$-primary.

For $x\in A$, write $(Q:x)=\{m\in M:xm\in Q\}$. If $xM\subseteq Q$, then $(Q:x)=M$. Otherwise multiplication by $x$ embeds $M/(Q:x)$ as a nonzero submodule of $M/Q$, so $(Q:x)$ is $\mathfrak{p}$-primary by the same argument. If $x\notin\mathfrak{p}$, multiplication by $x$ on $M/Q$ is injective, giving $(Q:x)=Q$.
